\documentclass[12pt]{article}
\usepackage[utf8]{inputenc}
\usepackage[T1]{fontenc}
\usepackage{hyperref}
\title{West Flemish Language Resources for K-12 Curriculum Development}
\author{Compiled by Claude Sonnet 5 (claude-sonnet-5) for Andr\'as Kornai}
\date{}
\begin{document}
\maketitle

\section{Introduction}
West Flemish (ISO 639-3 vls, Glottocode vlaa1240) is a Low Franconian variety spoken in West Flanders (Belgium), Zeeuws-Vlaanderen (Netherlands) and French Flanders. Speaker estimates vary by source, from about 1.05 million in Belgian West Flanders alone to about 1.6 million overall. Many linguists treat it as a group of Dutch dialects, while France, UNESCO, Glottolog and Ethnologue treat it as a distinct language. These searches found little on school materials in West Flemish.

\section{Primary Educational Resources}
In French Flanders, the Cassel-based Regional Flemish Language Institute (ANVT) is reported to publish books and teach the language to locals. I did not locate its own website in these searches.
\begin{itemize}
\item Brussels Times article on efforts to keep the Flemish dialect alive in France: \url{https://www.brusselstimes.com/world/1809587/the-corner-of-france-fighting-to-keep-its-flemish-dialect-alive}
\end{itemize}

\section{Linguistic and Reference Materials}
\begin{itemize}
\item West Flemish, Wikipedia: \url{https://en.wikipedia.org/wiki/West_Flemish}
\item Flemish dialects, Wikipedia: \url{https://en.wikipedia.org/wiki/Flemish_dialects}
\item Research portal record related to the Dictionary of Flemish Dialects (WVD, Ghent University, 1972--2018; online as e-WVD): \url{https://www.researchportal.be/en/node/942425}
\end{itemize}

\section{Religious and Community Resources}
No scripture or community organisation was found in the searches beyond the ANVT mention above.

\section{Digital and Audio Resources}
\begin{itemize}
\item Localazy language hub entry for vls (software localisation): \url{https://hub.localazy.com/en/languages/vls-west-flemish}
\item OpenStreetMap wiki, name:vls key (map labels in West Flemish): \url{https://wiki.openstreetmap.org/wiki/Key:name:vls}
\end{itemize}

\section{Curriculum Development Context}
The searches did not find West Flemish textbooks or classroom programmes. The main reference base is dialectological, in particular the WVD database. A curriculum would have to cope with the lack of a standardised orthography and with the status question of language versus dialect. A Wikipedia edition in vls may exist, but I did not confirm it from a page.

\end{document}